
Comme @@M0@@ par 1.7, @@M1@@ est @@M2@@-r\'egulier par 6.2\,; donc @@M3@@ pour @@M4@@ et @@M5@@ par 1.4 (i). La suite exacte @@M6@@ montre alors que @@M7@@ est @@M8@@-r\'egulier\,; d'o\`u (i).

L'assertion (i) implique que @@M9@@\,; d'o\`u (ii). D'apr\`es 2.10 et 1.9,
@@M10@@
donc, @@M11@@. Par r\'ecurrence sur @@M12@@, on peut supposer que @@M13@@ est un @@M14@@-faisceau\,; d'o\`u (iii) gr\^ace \`a 1.6.2.

\textbf{Lemme 6.5.} \emph{Sous les conditions 6.1, on suppose @@M15@@ sans cycle premier associ\'e de dimension @@M16@@.}
\begin{enumerate}[label=(\roman*)]
\item \emph{Si @@M17@@, alors @@M18@@.}
\item \emph{@@M19@@ pour @@M20@@.}
\item \emph{Soient une suite exacte @@M21@@ et un entier @@M22@@ tels que @@M23@@\,; alors @@M24@@ et @@M25@@.}
\end{enumerate}
En effet, soit @@M26@@, @@M27@@; alors @@M28@@ d\'efinit un sous-faisceau non nul @@M29@@ de @@M30@@, donc @@M31@@ est @@M32@@. Prenons un diagramme commutatif \`a lignes et colonnes exactes
@@M33@@
